% Original background stub (preserved for reference):
% Background. test

\section{Grid and Carbon Fundamentals}

\subsection{How Carbon Intensity is Calculated}

The environmental impact of electricity consumption is quantified through carbon intensity, measured in grams of CO\textsubscript{2} equivalent per kilowatt-hour (g$\cdot$CO\textsubscript{2}eq/kWh).
This metric captures the amount of greenhouse gas emitted per unit of electricity generated, accounting for multiple gases through their CO\textsubscript{2} equivalent (CO\textsubscript{2}eq) warming potential.
Each energy source carries a distinct carbon factor, defined as the amount of CO\textsubscript{2}eq emitted per kWh of electricity it produces.
Carbon factors are calculated per energy source by the Intergovernmental Panel on Climate Change (IPCC) and represent the direct emissions from the operational generation of electricity \cite{ipcc_carbon_factors_2014}.

Two common signals are used to quantify the carbon intensity of grid electricity \cite{sukprasert_limitations_2024}.
The \textit{average carbon intensity} is the weighted average of the carbon factors of all generators currently supplying electricity to the grid, where each source is weighted by its share of total generation.
The \textit{marginal carbon intensity} represents the carbon factor of the specific generator that responds to the last incremental unit of demand.
Because not all generators respond equally to changes in demand, only the marginal generator is credited with the emissions associated with added consumption.
Although marginal intensity more precisely captures the impact of an additional unit of demand, it is highly volatile and difficult to observe in real time.
The average intensity signal, by contrast, is stable and directly observable from publicly reported generation data, and is the signal required for Scope 2 greenhouse gas reporting under the GHG Protocol \cite{ghg_protocol}.
Consistent with prior work \cite{maji_carboncast_2022,yan_ensembleci_2025,zhang_gnn-based_2023}, this thesis focuses on average carbon intensity.

The average carbon intensity of a region at a given time is computed as:
\begin{equation}
    CI_t = \frac{\sum_{i} E_{i,t} \cdot CF_i}{\sum_{i} E_{i,t}}
\end{equation}
where $E_{i,t}$ is the energy generated by source $i$ at time $t$ (in kWh) and $CF_i$ is the carbon factor of source $i$ (in g$\cdot$CO\textsubscript{2}eq/kWh).
The resulting carbon intensity reflects the instantaneous emissions profile of the grid based on its current generation mix.

\subsection{Energy Source Mix}

An electricity grid draws from a portfolio of generation sources that differ in their carbon factors, dispatchability, and cost.
These sources fall into three broad categories: fossil fuels, renewables, and low-carbon sources.

Fossil fuels, primarily coal and natural gas, carry high carbon factors.
Coal is among the most carbon-intensive sources, emitting approximately 820~g$\cdot$CO\textsubscript{2}eq/kWh, while natural gas emits roughly 490~g$\cdot$CO\textsubscript{2}eq/kWh \cite{ipcc_carbon_factors_2014}.
Despite their environmental cost, fossil fuels are highly dispatchable: operators can increase or decrease their output on demand, making them a reliable tool for balancing supply and demand.

Renewable sources, including wind, solar, and hydropower, have near-zero direct carbon factors.
Solar photovoltaic systems emit approximately 41~g$\cdot$CO\textsubscript{2}eq/kWh and onshore wind approximately 11~g$\cdot$CO\textsubscript{2}eq/kWh over their lifecycle \cite{ipcc_carbon_factors_2014}.
Unlike fossil fuels, however, wind and solar are non-dispatchable: their output depends entirely on prevailing weather conditions and cannot be controlled by grid operators.
Hydropower is partially dispatchable, as reservoir levels can be managed to modulate output, but is subject to longer-term seasonal constraints.

Low-carbon sources include nuclear power and battery storage.
Nuclear power produces approximately 12~g$\cdot$CO\textsubscript{2}eq/kWh and provides stable, dispatchable baseload generation.
Battery storage systems have near-zero direct emissions and serve primarily as a buffer, storing excess generation from renewables and discharging during periods of high demand.

Each regional grid is managed by a balancing authority, which is responsible for continuously matching electricity supply with demand in real time.
The carbon intensity of a grid at any moment is determined by which generators the balancing authority has dispatched to meet current demand.
Regions with a high share of renewables and nuclear in their generation mix, such as France (approximately 70\% nuclear) or Sweden (dominated by hydro and nuclear), have persistently low carbon intensity.
Regions that rely heavily on coal and gas, such as parts of the United States Midwest or Poland, maintain high carbon intensity.

% TODO: Insert Figure -- spatial variability bar chart (carbon intensity by region, sorted by yearly average)

\subsection{Temporal and Spatial Variability}

Carbon intensity is not static: it varies continuously as the generation mix shifts in response to demand, weather, and grid conditions.
This variability operates across two dimensions, temporal and spatial, and understanding both is essential for effective carbon-aware scheduling.

\textit{Temporal variability} arises from the interaction of demand patterns and renewable generation cycles.
Electricity demand follows predictable diurnal rhythms tied to human behavior: demand is lowest overnight, rises sharply in the morning as households and businesses come online, and peaks in the early evening when both commercial and residential activity are at their highest.
During periods of high demand, grid operators dispatch additional generators to maintain supply, often relying on carbon-intensive peaking units such as natural gas turbines, which temporarily elevates carbon intensity.
Renewable generation adds a second layer of temporal variability.
Solar output follows a strict daily cycle, peaking around midday and dropping to zero at night.
In regions with high solar penetration, such as California, this creates a pronounced mid-day dip in carbon intensity as solar displaces fossil fuel generation, followed by a rise in the evening as solar output falls and demand remains high.
Wind generation is less predictable, varying over hours and days depending on local and regional weather patterns.
Seasonal effects compound these daily cycles: solar irradiance is lower in winter, heating demand rises in cold months, and wind patterns shift between seasons, each affecting the generation mix and the resulting carbon intensity.
Weather also affects demand directly: temperature extremes drive energy consumption for heating and cooling, which is frequently met by carbon-intensive peaking units, raising carbon intensity during heat waves and cold snaps.

\textit{Spatial variability} reflects the fundamental differences in generation infrastructure across regions.
Carbon intensity can vary by more than 10$\times$ between the cleanest and most carbon-intensive grids worldwide.
For instance, France's nuclear-heavy grid maintains an average carbon intensity near 60~g$\cdot$CO\textsubscript{2}eq/kWh, while coal-dependent regions such as Montana or Poland average above 600~g$\cdot$CO\textsubscript{2}eq/kWh.
Even within a single country, carbon intensity varies significantly between regions: California's grid, with a high share of solar and wind, is substantially cleaner than the coal-heavy Midwest.

% TODO: Insert Figure -- temporal variability line plot (hourly CI over a 7-day window for representative regions)

\subsection{Intermittency of Renewable Energy}

Wind and solar energy are inherently intermittent: they generate electricity only when weather conditions permit and their output cannot be controlled by grid operators.
Unlike coal or gas plants, which can be ramped up or shut down on demand, renewable generators are non-dispatchable.
A sudden drop in wind speed or an overcast sky can cause renewable output to fall within minutes, forcing the grid to compensate with fast-ramping fossil fuel generators.
This non-dispatchability is a fundamental physical constraint of renewable generation.

As electricity grids integrate higher shares of renewables to reduce their carbon footprint, this intermittency becomes an increasingly important source of carbon intensity variability.
Grids with high renewable penetration, such as Denmark (dominated by wind) or Texas (significant wind capacity), exhibit large and rapid swings in carbon intensity that are difficult to predict.
The variability in these regions is substantially higher than in grids anchored by stable baseload sources such as nuclear or hydro.

This layered uncertainty, from weather conditions to renewable output to carbon intensity, makes precise long-horizon carbon intensity forecasting inherently difficult.
A forecaster must anticipate not only how demand will evolve but also how wind speeds and solar irradiance will change, and how the grid will respond to those conditions.
This motivates the need for forecasting approaches that are adaptive to changing grid conditions rather than reliant solely on historical patterns learned from large datasets.

\section{Carbon-Aware Scheduling}

\subsection{Types of Electrical Loads}

An electrical load is considered flexible when it can adapt its consumption pattern in response to the availability of low-carbon electricity.
Flexibility has two independent dimensions: temporal and spatial.

\textit{Temporal flexibility} is the ability of a load to be postponed, extended, or shortened relative to its nominal schedule.
Temporally flexible loads fall into two categories based on whether their execution can be interrupted.
A \textit{continuous} load must run uninterrupted from start to finish once it begins: stopping and restarting would either damage the process or produce an incorrect result.
Examples include aluminum smelting, AI model training, and batch scientific computing jobs.
An \textit{interruptible} load can be paused and resumed at any point without affecting its outcome, provided it completes the required total amount of work within a deadline.
EV charging is a canonical example: a vehicle must receive a target charge by a certain time, but the individual charging sessions can be distributed across that window in any pattern.
Residential appliances such as dishwashers and washing machines are similarly interruptible.

\textit{Spatial flexibility} is the ability of a load to shift its execution to a different geographic location.
Spatially flexible loads can migrate across multiple sites, each connected to a different regional grid with a potentially different carbon intensity.
Web services and cloud batch workloads are common examples: a distributed cluster spanning multiple regions can redirect computation to whichever site currently has the lowest carbon intensity.
Like temporal flexibility, spatial flexibility can combine with either continuous or interruptible execution.
A web service is spatially flexible and continuous, since it must remain available at all times but can migrate with minimal disruption.
A batch job may be both temporally and spatially flexible, allowing it to wait for a cleaner time as well as migrate to a cleaner region.

\subsection{Carbon-Aware Optimization Strategies}

Carbon-aware optimizers exploit the temporal and spatial variability of carbon intensity to reduce the operational carbon footprint of flexible workloads.
Two principal strategies are used, corresponding to the two dimensions of load flexibility.

\textit{Temporal shifting} defers or interrupts a workload to align its execution with periods of lower carbon intensity within a deadline window.
For a continuous workload, this means delaying the start time until a low-carbon period is forecast to begin.
For an interruptible workload, the scheduler selects the lowest-carbon hours within the deadline window and distributes execution across them.
The amount of carbon savings achievable through temporal shifting depends on the magnitude of temporal variability in the local grid and the amount of slack available in the workload's deadline \cite{sukprasert_limitations_2024}.
Google's carbon-intelligent computing platform applies temporal shifting to defer batch datacenter workloads to hours when renewable energy supply is highest \cite{google_data_center_work_harder}.
Wiesner et al.\ demonstrate how temporal shifting can reduce cloud carbon emissions across different European regions \cite{wiesner_lets_2021}.

\textit{Spatial shifting} migrates a workload to a geographic region with lower current carbon intensity.
This strategy is applicable to workloads with spatial flexibility, such as distributed cloud services.
CASPER \cite{souza_casper_2023} formulates spatial shifting for distributed web services as a multi-objective optimization problem balancing carbon footprint and latency constraints, achieving up to 70\% carbon reduction.
Caribou \cite{gsteiger_caribou_2024} applies fine-grained geospatial shifting to serverless functions, enabling sub-request-level migration to cleaner regions.

The two strategies are complementary and can be combined for workloads that possess both temporal and spatial flexibility.
The degree of achievable savings depends on the variability of carbon intensity across time and space, as well as the flexibility of the workload.
Sukprasert et al.\ provide a comprehensive analysis of the practical upper bounds of both strategies across 123 cloud regions, showing that while carbon savings are meaningful, they are bounded and diminish as grids become cleaner \cite{sukprasert_limitations_2024}.

\subsection{Carbon Forecasting for Scheduling}

Carbon-aware optimization requires planning ahead.
To schedule a workload at the optimal time or location, an optimizer must know when and where carbon intensity will be lowest before the workload begins.
This planning requires forecasts: without a look-ahead window of future carbon intensity, a scheduler cannot identify the best execution window and must resort to reacting to current conditions rather than anticipating future ones \cite{grid_observability}.
Carbon intensity forecasting is therefore a critical enabler of carbon-aware optimization.

Effective forecasting must provide predictions over the full scheduling horizon of the workload, which can range from a few hours for short batch jobs to multiple days for long-running AI training workloads with large deadline slack.
Current state-of-the-art forecasting approaches build complex neural network models to capture the spatiotemporal dynamics of the grid's energy source mix.
CarbonCast \cite{maji_carboncast_2022} uses a two-tier hierarchical approach: a first tier of neural networks generates production forecasts for each individual energy source, and a second tier combines these with historical carbon data and weather forecasts using a CNN-LSTM architecture to produce multi-day carbon intensity forecasts up to 96 hours ahead.
EnsembleCI \cite{yan_ensembleci_2025} extends this approach by combining multiple sub-learner models through ensemble learning, selecting the best predictor per region to improve accuracy and reduce variability across regions.
Zhang and Wang \cite{zhang_gnn-based_2023} apply graph neural networks to model the spatial dependencies between cross-border grids in Europe, capturing how electricity exchanges between neighboring countries affect each region's carbon intensity.

These models achieve strong accuracy results, with CarbonCast reporting a MAPE of 4.80--13.93\% across multiple regions.
However, they share two important limitations.
First, they require large volumes of historical energy, weather, and carbon data to train effectively, often spanning multiple years, and depend on significant computational resources.
This makes them difficult to deploy in regions with limited data availability and costly to retrain when grid conditions change.
Second, and more fundamentally, their effectiveness for carbon-aware scheduling has not been evaluated.
These models are optimized to minimize forecasting error metrics such as MAPE, but whether lower MAPE actually translates to better scheduling outcomes, and what property of a forecast is most consequential for scheduling, remains an open question that this thesis addresses.
